% Annotated listing of the reference implementation khinchin_ref.py
% (Theorem thm:khinchin), \input from zetaF.tex.  The line ranges below refer
% to khinchin_ref.py and must be updated if that file changes.

\section{Reference implementation for Khinchin's constant}\label{sec:khpseudo}

\lstdefinestyle{ref}{language=Python, basicstyle=\ttfamily\scriptsize, columns=fixed, basewidth=0.5em, keepspaces=true,
  numbers=left, numberstyle=\tiny, numbersep=6pt, xleftmargin=1.5em, showstringspaces=false,
  keywordstyle=\ttfamily, commentstyle=\itshape, aboveskip=0.6\baselineskip, belowskip=0.6\baselineskip}
\newcommand{\reflisting}[2]{\lstinputlisting[style=ref, firstline=#1, lastline=#2, firstnumber=#1]{khinchin_ref.py}}

The Python program \texttt{khinchin\_ref.py} in the repository evaluates Theorem~\ref{thm:khinchin} with the
bounds of Section~\ref{sec:khbounds}. We list it in full, with annotations; the line numbers are those of the
file. It uses python-flint: \texttt{fmpq} and \texttt{fmpq\_mat} are exact rationals and rational matrices, and
\texttt{arb} is a real ball. Quantities computed from $u_0$, $k$, $n$, $j$, $c_0$, $c$ and $\delta$ alone are exact
rationals, and so are all matrices. Every other quantity is a ball at working precision $p + 30$ bits, and
\texttt{pm(x, r)} adds $r$ to the radius of $x$. The program is written for clarity, not speed; at $p = 5000$ it
takes a few seconds and returns a ball of radius $6\cdot10^{-1507}$ containing $K$.

\emph{Parameters and helpers.} The split point $u_0 = \frac1{16}$, the number $r_0 = 16$ of bits of the first
bit-burst point and the growth factor $g = 4$ are global. A line \texttt{N = least(n0, cond)} sets $N$ to the least
$N \ge n_0$ for which an upper bound of a tail bound $R(N)$ is at most $\varepsilon = 2^{-p-10}$. The black boxes are
$\exp$, $\sqrt{\cdot}$, $\erfc$, $\pi$, $\log 2$, $\Li_2(-\frac12)$ and $\Eip(y) = \operatorname{Ei}(y) - \gamma - \log y$.
\reflisting{12}{37}

\emph{Binary splitting.} All series are summed by one primitive. If $v_{n+1} = T_nv_n$ with
$T_n \in \QQ^{d\times d}$ and $W_n \in \QQ^{e\times d}$, then $v_b = Pv_a$ and $\sum_{n=a}^{b-1}W_nv_n = Sv_a$. An
efficient implementation writes $T_n = \tilde T_n/q_n$ with integer $\tilde T_n$ and $q_n$, uses the triangular or
sparse structure of the $T_n$, and applies partial products to the state instead of forming $S$ exactly
(Section~\ref{sec:khimpl}).
\reflisting{39}{47}

\emph{The $k$-terms.} $I_k = e^{-x}\sum_{n=1}^{\infty}P_n/n^2$ is summed with the state
$(P_n, Q_n) = ((4/K)^ne_n(x), (4u_0)^n/n!)$, starting from $(P_1, Q_1) = (q, 4u_0)$, where $K = k^2$, $x = Ku_0$ and
$q = 4/K$; $R(N)$ is the tail bound of Section~\ref{sec:khbounds}.
\reflisting{49}{68}

\emph{The main term.} Since $\pi_n(y) = e^{-y}\sum_{i=n}^{\infty}y^i/i!$, exchanging the order of summation gives
$M = e^{-u_0}A - e^{-4u_0}B - \frac12C$ with $A = \sum_{i=1}^{\infty}w_i$, $B = \sum_{i=1}^{\infty}\hat w_i$ and
$C = \sum_{i=1}^{\infty}v_i/i^2$, where
\[
  v_i = \frac{(4u_0)^i}{i!}, \qquad w_i = \frac{u_0^i}{i!}\sum_{n=1}^{i}\frac{4^n}{n^2}, \qquad
  \hat w_i = v_i\sum_{n=1}^{i}\frac1{n^2},
\]
and $v_1 = w_1 = \hat w_1 = 4u_0$; the state is $(v_i, w_i, \hat w_i)$.
\reflisting{70}{82}

\emph{Starting values.} The dual state at $t = \frac32$ is
\[
  V = \bigl(\nu_{3/2},\ Y_{3/2},\ Y_{1/2},\ C_{3/2},\ \operatorname{Forcing}(E)\bigr),
\]
where the last three components are $E\,u_0^{-1/2}(\Eip(4u_0), e^{4u_0}, 1)$. At a rational point $c_0$, with $D = D_{-1/2}(c_0)$ and
$\hat E = e^{-c_0/u_0}u_0^{-1/2}$, the recurrence \eqref{eq:Drec} gives $D_{n-1/2} = \alpha_nD + \tilde\beta_n\hat E$
with $\alpha_0 = 1$, $\tilde\beta_0 = 0$ and
\[
  \alpha_{n+1} = -\frac{c_0\alpha_n}{n+\frac12}, \qquad \tilde\beta_{n+1} = \frac{u_0^{n+1} - c_0\tilde\beta_n}{n+\frac12}.
\]
With $b_n = 4^n/n!$ the binary splitting state is $(b_n\alpha_n, b_n\tilde\beta_n, b_nu_0^n)$, and the six rows of
$W_n$ are the $n$-th terms of the $D$- and $\hat E$-parts of $\nu_{3/2} = \sum_{n=1}^{\infty}\frac{b_n}nD_{n-1/2}$,
$Y_{3/2} = \sum_{n=0}^{\infty}b_nD_{n-1/2}$ and $Y_{1/2} = \sum_{n=0}^{\infty}b_nD_{n+1/2}$; the terms $n = 0$ are added
separately (line~116). The two parts have size about $e^{4\sqrt{c_0} - c_0/u_0}$ and their sum about $e^{-c_0/u_0}$;
the cancellation costs no absolute precision, since both parts are smaller than $1$.
\reflisting{84}{120}

\emph{Taylor stages.} At a fixed point $c$, let $V_j$ be the state at $t = \frac32 + j$,
\[
  V_j = \bigl(\nu_t,\ Y_t,\ Y_{t-1},\ C_t,\ \Eip(4u_0)Eu_0^{1-t},\ e^{4u_0}Eu_0^{1-t},\ Eu_0^{1-t}\bigr).
\]
The integration-by-parts recurrences of Section~\ref{sec:dilog} read $V_{j+1} = M_jV_j$ with
\[
  M_j = \frac1c\begin{pmatrix}
  \eta_j & -1 & 0 & 1 & 1 & 0 & 0 \\
  0 & \eta_j & -4 & 0 & 0 & 1 & 0 \\
  0 & c & 0 & 0 & 0 & 0 & 0 \\
  0 & 0 & 0 & \eta_j & 0 & 0 & 1 \\
  0 & 0 & 0 & 0 & c/u_0 & 0 & 0 \\
  0 & 0 & 0 & 0 & 0 & c/u_0 & 0 \\
  0 & 0 & 0 & 0 & 0 & 0 & c/u_0
  \end{pmatrix}, \qquad \eta_j = j + \tfrac12.
\]
Since $\partial_cZ_t = -Z_{t+1}$ and $E(c+\delta) = E(c)e^{-\delta/u_0}$, every component satisfies
$V(c+\delta) = \sum_{j=0}^{\infty}\frac{(-\delta)^j}{j!}V_j(c)$, so $T_j = -\frac{\delta}{j+1}M_j$. The first four
components are summed this way and the last three are recomputed.
\reflisting{122}{142}

\emph{Dual terms.} This is the bit-burst algorithm. The point $c_m = \pi^2(m+\frac12)^2$ is truncated to
$r_0, gr_0, g^2r_0, \dots$ fractional bits, up to the number $r^*$ of bits required by the bound for $|J'|$ in
Section~\ref{sec:khbounds}, and $J_m = \sqrt\pi\,\nu_{3/2}(c_m)$.
\reflisting{144}{160}

\emph{Khinchin's constant.} The functions \texttt{tk} and \texttt{tm} are the bounds for the omitted $k$-terms and
dual terms of Section~\ref{sec:khbounds}, and the constant is assembled as in Theorem~\ref{thm:khinchin}. For
example, $p = 5000$ gives $k_{\max} = 235$ and $m_{\max} = 4$, and the dual term $m = 0$ uses the points with
$r = 16, 64, 256, 1024, 4096, 4951$ fractional bits, and $m = 4$ those with $r = 16, 64, 256, 391$.
\reflisting{162}{182}
